\documentclass[11pt, a4paper]{amsart}
\usepackage{amsmath, amssymb, amsthm, amsfonts}
\usepackage{mathtools}
\usepackage[margin=1in]{geometry}
\usepackage{hyperref}

\newtheorem{theorem}{Theorem}[section]
\newtheorem{lemma}[theorem]{Lemma}
\newtheorem{proposition}[theorem]{Proposition}
\newtheorem{corollary}[theorem]{Corollary}
\newtheorem{fact}[theorem]{Standard Fact}
\theoremstyle{definition}
\newtheorem{definition}[theorem]{Definition}
\newtheorem{remark}[theorem]{Remark}
\newtheorem{notation}[theorem]{Notation}
\newtheorem{hypothesis}[theorem]{Hypothesis}

\newcommand{\R}{\mathbb{R}}
\newcommand{\Q}{\mathbb{Q}}
\newcommand{\Z}{\mathbb{Z}}
\newcommand{\N}{\mathbb{N}}
\newcommand{\id}{\operatorname{id}}
\newcommand{\tr}{\operatorname{tr}}
\newcommand{\spec}{\operatorname{spec}}
\newcommand{\op}{\operatorname{op}}
\newcommand{\Sym}{\operatorname{Sym}}
\newcommand{\supp}{\operatorname{supp}}
\newcommand{\dsph}{d_{S^3}}
\newcommand{\Pihyd}{\Pi_{\mathrm{hyd}}}
\newcommand{\Lmacro}{\mathcal L_{\mathrm{macro}}}
\newcommand{\divS}{\operatorname{div}^{S^3}}
\newcommand{\LapS}{\Delta_{S^3}}
\newcommand{\graphLap}{\Delta^{\rm graph}_n}
\newcommand{\Projn}{\mathbb P_n}
\newcommand{\Vdf}{\mathcal V_n^{\mathrm{df}}}

\title{Cascade--Classical Correspondence VII:\\
  Hydrodynamic Projection and Scaling Limits}
\author{}
\date{}

\begin{document}
\maketitle

\begin{abstract}
On the $D_4$ refinement branch of
\cite{MetricProjection, SchlafliConvergence},
and using the finite-level closure dynamics of
\cite{ClosureDynamics}, we construct an explicit hydrodynamic
projector $\Pihyd^{(n)} \colon X_n \to \mathcal U_n \times
\mathcal P_n$ from the P4 state space to a velocity--pressure
space on the finite-level geometry. The paper has two
structurally distinct layers:

\begin{enumerate}
\item \textbf{Unconditional finite-level layer.} For every
  level $n$, we state and give proof outlines for:
\begin{itemize}
\item[\textbf{T1}] (Theorem~\ref{thm:pihyd-bounded})
  $\Pihyd^{(n)}$ is well-defined and bounded in explicit norms
  with $\|u_n[\Phi]\|_{\mathcal U_n} + \|p_n[\Phi]\|_{\mathcal
  P_n} \leq C_n \|\Phi\|_{X_n}$, and its divergence admits the
  printed decomposition $\divS u_n[\Phi] =
  \mathcal C_n[\Phi] = A_n[\Phi] + B_n[\Phi] + E_n[\Phi]$ into
  three explicit named terms, with the parallel-transport
  correction $E_n$ \emph{defined} to make the identity hold
  (not a substantive compatibility theorem between the P4
  discrete divergence and the continuum divergence).
\item[\textbf{Prop~\ref{prop:balance-laws}}] Finite-level
  mass / momentum / energy balance laws along the P4 flow.
\item[\textbf{Prop~\ref{prop:leray}}] Finite-level
  Leray--Helmholtz decomposition $\mathcal V_n = \Vdf \oplus
  \nabla \mathcal Q_n$ with bounded projector $\Projn$.
\item[\textbf{T2}] (Theorem~\ref{thm:projected-dynamics}) The
  projected state $(u_n, p_n) := \Pihyd^{(n)} \Phi_n(t)$
  satisfies, in macroscopic time $\tau = \lambda_n t$,
  \[
    \partial_{\tau, n} \Pihyd^{(n)} \Phi_n
      = \Lmacro^{(n)}\bigl(\Pihyd^{(n)} \Phi_n\bigr)
      + D_n(\Phi_n),
  \]
  where $D_n(\Phi) = (D_n^{\rm vel}(\Phi), D_n^{\rm press}(\Phi))
  \in L^2(S^3; TS^3) \times L^2(S^3)$ is a product-space object
  whose velocity component decomposes as $D_n^{\rm vel} =
  D_n^{\rm comm} + D_n^{\rm geom} + D_n^{\rm div} + D_n^{\rm
  scale}$, with all five named components ($D_n^{\rm avg} \equiv
  0$ for the time-independent kernel) printed by closed formula,
  and $L^2$-norm rates $O(h_n^{1/2})$ stated in
  Proposition~\ref{prop:defect-bounds} at the outline level,
  with concrete constants and full proofs deferred.
\end{itemize}
\item \textbf{Conditional program (under
  Hypothesis~\ref{hyp:scaling})}:
\begin{itemize}
\item[\textbf{T3}] (Theorem~\ref{thm:ns-limit}) Under the
  explicit scaling / regularity / compactness hypothesis package
  $\mathbf H_{\rm hyd}$ stated before use, the total defect
  $\delta_n(T) \to 0$ in the printed norm, the projected fields
  are precompact in appropriate function spaces, and every limit
  point $(u, p)$ solves incompressible Navier--Stokes on the
  printed limit geometry in the exact weak/strong sense fixed in
  the theorem, with viscosity $\nu := \lim_n \nu_n$.
\end{itemize}
\end{enumerate}

The unconditional layer uses only P4, P5, and the P6
\emph{unconditional extrinsic layer}
\cite[Thms.~4.2, 5.1, Prop.~6.1, Prop.~7.2]{SchlafliConvergence}.
The conditional layer may additionally use P6's conditional
program (intrinsic metric, form consistency, spectral
convergence) under its Hypothesis $6.4$, and invokes standard
PDE compactness / NS weak-solution theory.

We make \emph{no} unconditional Navier--Stokes, regularity,
uniqueness, blow-up, or smoothness claim. We make \emph{no}
Euler claim in the main theorem (Euler is treated, if at all,
as a separate inviscid corollary with its own hypothesis
bundle). We do not cite \texttt{cascade-ns-proof.md},
\texttt{ns-formal.tex}, or \texttt{foundations.tex} as
theorem-grade proof support.
\end{abstract}

\tableofcontents

%=====================================================================
\section{Introduction and Source Audit}
\label{sec:intro}
%=====================================================================

\subsection{What P7 is and is not}

P7 is the \emph{hydrodynamic bridge} in the cascade
infrastructure: it takes the finite-level closure dynamics of
P4 \cite{ClosureDynamics}, the metric-side projections of P5
\cite{MetricProjection}, and the unconditional geometric
convergence of P6 \cite{SchlafliConvergence}, and produces a
velocity--pressure observable $\Pihyd^{(n)}$ together with a
projected evolution equation with fully explicit defect terms.
Under a stated scaling and regularity hypothesis, it derives
incompressible Navier--Stokes on the printed limit geometry as
the conditional scaling limit of the defect-corrected projected
dynamics.

P7 is \emph{not}:
(a) an unconditional Navier--Stokes theorem;
(b) a regularity, uniqueness, smoothness, or blow-up result;
(c) a millennium-level NS paper; those belong to the downstream
NS paper;
(d) a global well-posedness or bounded-energy statement beyond
the printed hypothesis bundle;
(e) an Euler theorem; Euler is an optional separate corollary
with its own explicit inviscid hypotheses;
(f) a rederivation of unproven claims from
\texttt{cascade-ns-proof.md} or \texttt{ns-formal.tex}; those
are deprecated and not cited.

The split between unconditional and conditional layers is
rigorously enforced: the unconditional theorems
(Theorem~\ref{thm:pihyd-bounded},
Proposition~\ref{prop:balance-laws},
Proposition~\ref{prop:leray},
Theorem~\ref{thm:projected-dynamics},
Proposition~\ref{prop:defect-bounds}) use only the P4, P5, and
P6-extrinsic-layer theorem-grade inputs; the conditional theorem
(Theorem~\ref{thm:ns-limit}) explicitly assumes
Hypothesis~\ref{hyp:scaling}, which includes any conditional P6
hypotheses used.

\subsection{Source audit}

\textbf{Cascade-internal:}
\begin{itemize}
\item P4~\cite{ClosureDynamics}: finite-level state spaces $X_n
  = X_n^{0, D_4} \oplus X_n^{1, D_4}$
  \cite[Def.~2.1]{ClosureDynamics}, closure generator $L_n$, flow
  $\dot\Phi_n = -L_n \Phi_n$, semigroup $e^{-tL_n}$
  \cite[Thm.~4.2]{ClosureDynamics}, energy dissipation
  \cite[Prop.~5.1]{ClosureDynamics}.
\item P5~\cite{MetricProjection}: $D_4$ coordinates
  $x_v^{(n)}$, finite-level metric tensor package.
\item P6~\cite{SchlafliConvergence}: \emph{unconditional}
  extrinsic metric $d_n^{\rm ext}$, mesh scale $h_n$,
  $\varepsilon$-net rate $O(h_n)$, GH convergence to $(S^3,
  \dsph)$ (see P6 \cite[Thm.~4.2, Thm.~5.1]{SchlafliConvergence}),
  doubling / packing
  (Prop~6.1 of P6), $V_0$-orbit isotropy constant $8/3$
  (Prop~7.2 of P6).
\item P6 conditional program (intrinsic metric
  $d_n^{\rm int}$, Dirichlet-form consistency, spectral
  convergence) used only inside Theorem~\ref{thm:ns-limit}
  under Hypothesis~\ref{hyp:scaling}, which includes P6
  Hypothesis 6.4 (uniform $8/3$ approximate isotropy at every
  $V_n$ vertex).
\end{itemize}

\textbf{Classical (standard mathematics):}
\begin{itemize}
\item Finite-dimensional linear algebra, bounded-operator norms
  \cite{ReedSimonI}.
\item Discrete exterior calculus / discrete divergence on
  simplicial complexes \cite{Chung}.
\item Leray--Helmholtz decomposition on finite-dimensional
  inner-product spaces \cite{ArnoldFalkWinther}.
\item Compactness tools (Aubin--Lions--Simon, Banach--Alaoglu)
  \cite{SimonCompactness, BrezisFA}.
\item Navier--Stokes weak-solution theory on compact Riemannian
  manifolds \cite{TemamNS, LionsMathTopics}.
\item Round-sphere analysis and function spaces on $(S^3,
  g_{\rm round})$ \cite{BerardSpectralGeometry, TaylorPDE1}.
\end{itemize}

\textbf{Explicitly excluded:}
\texttt{cascade-ns-proof.md}, \texttt{ns-formal.tex},
\texttt{cascade-foundations.md}, \texttt{foundations.tex}.

\subsection{Organization}

\S\ref{sec:setup} recalls the P4/P5/P6 inputs.
\S\ref{sec:kernel} defines the averaging kernel $K_{n,\ell_n}$,
coarse-graining scale $\ell_n$, and the admissible regime.
\S\ref{sec:observables} defines the current / momentum
observable $J_n$, velocity $u_n$, pressure $p_n$, and the
hydrodynamic projector $\Pihyd^{(n)}$.
\S\ref{sec:bounded} states T1 (Theorem~\ref{thm:pihyd-bounded})
and gives its proof outline.
\S\ref{sec:balance} states and outlines the balance laws
(Proposition~\ref{prop:balance-laws}).
\S\ref{sec:leray} states and outlines the Leray--Helmholtz
decomposition (Proposition~\ref{prop:leray}).
\S\ref{sec:projected} states T2
(Theorem~\ref{thm:projected-dynamics}) with the full defect
decomposition and proof outline.
\S\ref{sec:defect-bounds} states and outlines proofs of the
defect estimates (Proposition~\ref{prop:defect-bounds}).
\S\ref{sec:scaling-limit} states T3
(Theorem~\ref{thm:ns-limit}) under
Hypothesis~\ref{hyp:scaling} with proof outline.
\S\ref{sec:audit}, \S\ref{sec:handoff}, \S\ref{sec:notation}
are self-audit, citation contract, and notation audit.

%=====================================================================
\section{Setup: P4 State, P5/P6 Geometry}
\label{sec:setup}
%=====================================================================

\begin{notation}[P4 state space and closure flow]
\label{not:P4}
Let $X_n := X_n^{0, D_4} \oplus X_n^{1, D_4}$ be the
finite-level $D_4$-branch state space of
\cite[Def.~2.1]{ClosureDynamics}; let $L_n \in B(X_n)$ be the
bounded self-adjoint closure generator of
\cite[Thm.~3.1]{ClosureDynamics}; let $\Phi_n(t) := e^{-t L_n}
\Phi_n(0) \in X_n$ denote the closure flow for initial data
$\Phi_n(0) \in X_n$.
\end{notation}

\begin{notation}[Geometric inputs from P5/P6 unconditional layer]
\label{not:geom}
Let $V_n := V_n^{D_4}$ be the level-$n$ $D_4$-branch vertex set
(imported via P5~\cite[Notation~2.1]{MetricProjection}); let
$x_v^{(n)} \in \R^4$ be
P5's canonical coordinate \cite[Notation~2.1]{MetricProjection};
let $\iota_n \colon V_n \to S^3$ be the iterated-spherical-midpoint
map of P6 \cite[Def.~3.2]{SchlafliConvergence}; let $d_n^{\rm
ext}(u, v) := \dsph(\iota_n(u), \iota_n(v))$ be the
\emph{unconditional} extrinsic metric. Let $h_n$ denote the mesh
scale \cite[Def.~3.4]{SchlafliConvergence}, $h_n \leq
\alpha^n \pi/2$ for any $\alpha > 1/\sqrt 3$, hence $h_n \to 0$
geometrically. Let $|V_n|$ satisfy $c_V h_n^{-3} \leq |V_n| \leq
C_V h_n^{-3}$ \cite[Lem.~6.2]{SchlafliConvergence}.
\end{notation}

%=====================================================================
\section{Averaging Kernel and Coarse-Graining Scale}
\label{sec:kernel}
%=====================================================================

\begin{definition}[Coarse-graining scale $\ell_n$]
\label{def:ell-n}
The \emph{hydrodynamic coarse-graining lengthscale} $\ell_n$ is
a printed sequence with $\ell_n \in (2 C_1 h_n, \pi/8)$ for all
$n$ large, where $C_1$ is the P6 $\varepsilon$-net constant. The
admissible regime is $h_n / \ell_n \to 0$ as $n \to \infty$
(i.e.\ $\ell_n$ is asymptotically much larger than the mesh).
An explicit choice is $\ell_n := h_n^{1/2}$, so that $h_n /
\ell_n = h_n^{1/2} \to 0$ and $\ell_n \to 0$.
\end{definition}

\begin{definition}[Macroscopic differential operators]
\label{def:macro-ops}
The averaging produces smooth vector / scalar fields on the
round sphere $(S^3, g_{\rm round})$. For these, we use the
standard Riemannian operators:
\begin{itemize}
\item $\nabla \colon C^\infty(S^3) \to C^\infty(S^3; TS^3)$ is
  the Levi-Civita gradient of $g_{\rm round}$.
\item $\operatorname{div}^{S^3} \colon C^\infty(S^3; TS^3) \to
  C^\infty(S^3)$ is the Riemannian divergence, adjoint of
  $-\nabla$ under the $L^2(S^3, \mu_\infty)$ inner product.
\item $\Delta_{S^3} := \operatorname{div}^{S^3} \circ \nabla$ is
  the (positive) Laplace--Beltrami operator on $C^\infty(S^3)$,
  extended to vector fields by the connection Laplacian.
\end{itemize}
These are defined entirely on the fixed Riemannian manifold
$(S^3, g_{\rm round})$ and do NOT depend on any finite-level
discretization. They are used \emph{unconditionally} throughout
this paper.

For the P4 discrete divergence on oriented edges, we reserve the
notation $\partial^\ast \colon X_n^{1, D_4} \to X_n^{0, D_4}$
from \cite[Def.~3.3]{ClosureDynamics}. The discrete graph
Laplacian of P6 (which depends on P6 Hypothesis~6.4 for
Laplace--Beltrami convergence) is called $\Delta^{\rm
graph}_n$ and is used ONLY in the conditional program
\S\ref{sec:scaling-limit}.
\end{definition}

\begin{definition}[Averaging kernel $K_{n, \ell_n}$]
\label{def:kernel}
For $x \in S^3$ and $v \in V_n$, define
\[
  K_{n, \ell_n}(x, v) \;:=\;
    \frac{\chi(\dsph(\iota_n(v), x)/\ell_n)}{Z_{n, \ell_n}(x)},
\]
where $\chi \colon [0, \infty) \to [0, 1]$ is a fixed smooth
bump function supported in $[0, 1]$ with $\chi(0) = 1$ and
$\chi \equiv 0$ on $[1, \infty)$, and
\[
  Z_{n, \ell_n}(x) \;:=\; \sum_{v \in V_n}
    \chi(\dsph(\iota_n(v), x)/\ell_n)
\]
is the normalization (summing only over $v \in V_n$ with
$\iota_n(v)$ within spherical distance $\ell_n$ of $x$).
\end{definition}

\begin{lemma}[Kernel properties]
\label{lem:kernel-properties}
For every $n$ with $h_n \leq h_\ast$ and every $x \in S^3$:
\begin{enumerate}
\item[\textup{(i)}] \emph{Partition of unity in $v$:}
  $\sum_v K_{n, \ell_n}(x, v) = 1$, so for any vertex-indexed
  quantity $\Psi(v)$,
  $\sum_v K_{n, \ell_n}(x, v)\,\Psi(v)$ is a local weighted
  average.
\item[\textup{(ii)}] \emph{Normalization is bounded above and
  below:} there exist $c_K, C_K > 0$ (independent of $n$) with
  $c_K\,(\ell_n/h_n)^3 \leq Z_{n, \ell_n}(x) \leq C_K\,
  (\ell_n/h_n)^3$. Hence $K_{n, \ell_n}(x, v) \leq C_K^{-1}\,
  (h_n/\ell_n)^3$ for every $v$.
\item[\textup{(iii)}] \emph{Support:} $K_{n, \ell_n}(x, v) = 0$
  if $\dsph(\iota_n(v), x) > \ell_n$.
\item[\textup{(iv)}] \emph{Smoothness in $x$:} $K_{n, \ell_n}(x,
  v)$ is $C^1$ in $x$ (since $\chi$ is smooth and $Z_{n,
  \ell_n}(x)$ is bounded below). $|\nabla_x K_{n, \ell_n}(x, v)|
  \leq C_K'\,(h_n/\ell_n)^3 / \ell_n$.
\end{enumerate}
\end{lemma}

\begin{proof}
(i): $\sum_v K(x, v) = \sum_v \chi_{x, v} / Z(x) = Z(x)/Z(x)
= 1$ whenever $Z(x) > 0$, which holds by (ii).
(ii): $Z(x) = \sum_{v : \dsph(\iota_n(v), x) \leq \ell_n}
\chi(\cdot) \geq $ count of $\iota_n(V_n)$ points in
$B_\dsph(x, \ell_n/2)$ (where $\chi \geq \chi(1/2) > 0$). By
P6 Prop~6.1 (doubling, unconditional), this count is $\Theta(
(\ell_n/h_n)^3)$ since $\ell_n \geq 2 C_1 h_n$. Upper bound from
the covering direction. (iii): immediate from $\chi = 0$ on
$[1, \infty)$. (iv): direct differentiation of $K(x, v) =
\chi(\cdot)/Z(x)$ using $|\nabla_x \chi(\dsph(\iota_n(v), x)/
\ell_n)| \leq \|\chi'\|_\infty / \ell_n$ and $|\nabla_x Z(x)|
\leq $ sum of $|\nabla_x \chi_{x,v}| \leq \|\chi'\|_\infty
|V_n^{\ell_n}(x)|/\ell_n$, where $V_n^{\ell_n}(x)$ is the
support-intersection set of size $\Theta((\ell_n/h_n)^3)$.
Combining gives the claimed bound on the gradient.
\end{proof}

%=====================================================================
\section{Velocity, Pressure, and the Hydrodynamic Projector}
\label{sec:observables}
%=====================================================================

\begin{definition}[Current observable $J_n$]
\label{def:Jn}
Given $\Phi = (f, B) \in X_n^{0, D_4} \oplus X_n^{1, D_4}$, the
\emph{discrete current / momentum observable} at $v \in V_n$ is
the tangent-vector-valued quantity
\[
  J_n[\Phi](v) \;:=\; \sum_{w : \{v, w\} \in E(G_n^{D_4})}
    B(\vec{vw})\,\hat\tau_{v, w} \;\in\; T_{\iota_n(v)} S^3,
\]
where $B(\vec{vw})$ is the oriented-edge coefficient from the
P4 state's 1-form part on the oriented edge $\vec{vw}$, and
$\hat\tau_{v, w} := \pi_v(\iota_n(w) - \iota_n(v)) /
\|\pi_v(\iota_n(w) - \iota_n(v))\|$ is the unit tangent
direction along edge $\{v, w\}$ (with $\pi_v$ the tangential
projection onto $T_{\iota_n(v)} S^3$).
\end{definition}

\begin{definition}[Averaged velocity and pressure]
\label{def:up}
The \emph{averaged velocity field} is
\[
  u_n[\Phi](x) \;:=\; \sum_{v \in V_n} K_{n, \ell_n}(x, v) \cdot
    \mathrm{TP}_v^x\bigl(J_n[\Phi](v)\bigr)
    \;\in\; T_x S^3,
\]
where $\mathrm{TP}_v^x \colon T_{\iota_n(v)} S^3 \to T_x S^3$ is
the spherical tangent-parallel-transport along the short
geodesic from $\iota_n(v)$ to $x$ (well-defined for
$\dsph(\iota_n(v), x) < \pi$, which holds on the kernel support
$\ell_n < \pi/8$).

The \emph{averaged pressure} is
\[
  p_n[\Phi](x) \;:=\; \sum_{v \in V_n} K_{n, \ell_n}(x, v)
    \cdot \rho_n[\Phi](v) \;-\; \frac{1}{\mu_\infty(S^3)}
    \int_{S^3} \sum_{v \in V_n} K_{n, \ell_n}(y, v) \cdot
    \rho_n[\Phi](v) \, d\mu_\infty(y),
\]
(the zero-mean normalization removes the constant-mode
ambiguity in the pressure), where $\rho_n[\Phi](v) :=
\tr_{F^0}(f(v))$ (trace of the
vertex-fibre component, as in P5's source extraction
\cite[Def.~3.2]{MetricProjection}) is the scalar
energy-density observable.
\end{definition}

\begin{definition}[Velocity and pressure function spaces]
\label{def:UP}
The \emph{velocity space} is
$\mathcal U_n := \{u \colon S^3 \to TS^3 : u \in C^0(S^3;
TS^3),\, u(x) \in T_x S^3\}$, with norm
$\|u\|_{\mathcal U_n} := \sup_{x \in S^3} |u(x)|_{g_{\rm round}}$.

The \emph{pressure space} is $\mathcal P_n := C^0(S^3)$ with
sup norm.

The \emph{state space norm} on $X_n$ is the P4 direct-sum norm
$\|\Phi\|_{X_n}^2 := \|f\|_{X_n^0}^2 + \|B\|_{X_n^1}^2$.
\end{definition}

\begin{definition}[Hydrodynamic projector $\Pihyd^{(n)}$]
\label{def:Pihyd}
The \emph{hydrodynamic projector} is the map
\[
  \Pihyd^{(n)} \colon X_n \to \mathcal U_n \times \mathcal P_n,
  \qquad \Pihyd^{(n)}(\Phi) \;:=\; \bigl(u_n[\Phi],
  p_n[\Phi]\bigr).
\]
\end{definition}

%=====================================================================
\section{T1: Boundedness and Divergence Decomposition}
\label{sec:bounded}
%=====================================================================

\begin{theorem}[Bounded hydrodynamic averaging]
\label{thm:pihyd-bounded}
For every $n$ with $h_n \leq h_\ast$, the hydrodynamic
projector $\Pihyd^{(n)}$ is a bounded linear map with
\[
  \|u_n[\Phi]\|_{\mathcal U_n} + \|p_n[\Phi]\|_{\mathcal P_n}
    \;\leq\; C_\Pi^{(n)} \|\Phi\|_{X_n},
\]
with explicit constant $C_\Pi^{(n)} = C_{\rm vel}\,h_n^{-3/2}$,
where $C_{\rm vel}$ is the base-geometric constant derived in
the proof (see Step~\emph{Boundedness} below). The $h_n^{-3/2}$
scaling comes from the pointwise-to-$\ell^2$ Cauchy--Schwarz step
over $|V_n| = \Theta(h_n^{-3})$ vertices; we do NOT claim
$n$-uniform boundedness of $\Pihyd^{(n)}$ in the sup norm, only
the per-level bound stated here.

Moreover, $u_n[\Phi]$ admits the following \emph{divergence
decomposition} (we do not claim this as a compatibility theorem
between the P4 discrete divergence $\partial^\ast B$ and the
continuum divergence $\divS$; it is a printed split of the
mismatch into three named terms, one of which is the
parallel-transport defect $\mathcal E^{\rm TP}_v[\Phi](x)$
\emph{defined} to make the identity hold): for every $x \in S^3$,
\[
  (\divS u_n[\Phi])(x) \;=\; \mathcal C_n[\Phi](x)
    \;:=\; A_n[\Phi](x) + B_n[\Phi](x) + E_n[\Phi](x),
\]
where the three explicit terms are:
\begin{align*}
  A_n[\Phi](x)
    &\;:=\; \sum_{v \in V_n}
       \bigl\langle \nabla_x K_{n, \ell_n}(x, v),\,
       \mathrm{TP}_v^x J_n[\Phi](v)\bigr\rangle_{g_{\rm round}}
    && \text{\emph{(kernel-gradient term)}} \\
  B_n[\Phi](x)
    &\;:=\; -\sum_{v \in V_n} K_{n, \ell_n}(x, v) \cdot
       (\partial^\ast B)(v)
    && \text{\emph{(local-divergence scalar)}} \\
  E_n[\Phi](x)
    &\;:=\; \sum_{v \in V_n} K_{n, \ell_n}(x, v) \cdot
       \mathcal E^{\rm TP}_v[\Phi](x)
    && \text{\emph{(parallel-transport correction)}}
\end{align*}
All three terms $A_n, B_n, E_n$ are scalar-valued functions on
$S^3$ (so the identity $\divS u_n = A_n + B_n + E_n$ is typed
correctly). The P4 discrete divergence of
\cite[Def.~3.3]{ClosureDynamics} applied to the edge component
is $\partial^\ast B \colon V_n \to \R$. The \emph{exact}
parallel-transport correction is
\[
  \mathcal E^{\rm TP}_v[\Phi](x) \;:=\;
    \divS\bigl(\mathrm{TP}_v^x J_n[\Phi](v)\bigr)(x)
    + (\partial^\ast B)(v),
\]
i.e.\ the defect between the Riemannian divergence of the
transported tangent vector field $x \mapsto \mathrm{TP}_v^x
J_n[\Phi](v)$ and the discrete local divergence at $v$. This is
a scalar-valued function of $x$. All three terms $A_n, B_n,
E_n$ are printed by closed formula; no $O(\cdot)$ placeholder
appears.

On the \emph{divergence-free source subspace}
$X_n^{\rm df} := \{\Phi \in X_n : (\partial^\ast B)(v) = 0
\text{ for all } v \in V_n\}$, the local-divergence term
$B_n[\Phi] \equiv 0$, and $\mathcal C_n[\Phi] = A_n[\Phi] +
E_n[\Phi]$.
\end{theorem}

\begin{proof}[Proof outline]
\emph{Boundedness.} By Definition~\ref{def:up} and
Lemma~\ref{lem:kernel-properties}(ii)(iv):
\begin{align*}
  |u_n[\Phi](x)|_{g_{\rm round}}
    &\;\leq\; \sum_v K_{n, \ell_n}(x, v) \cdot
       |J_n[\Phi](v)|_{g_{\rm round}} \\
    &\;\leq\; (\max_v K_{n, \ell_n}(x, v)) \cdot
       |V_n^{\ell_n}(x)| \cdot \max_v |J_n[\Phi](v)| \\
    &\;\leq\; C_K^{-1}(h_n/\ell_n)^3 \cdot C_K'(\ell_n/h_n)^3
       \cdot \max_v |J_n[\Phi](v)|
    \;\leq\; C'' \max_v |J_n[\Phi](v)|.
\end{align*}
For $|J_n[\Phi](v)|$: by Def~\ref{def:Jn}, $J_n[\Phi](v) =
\sum_{w \sim v} B(\vec{vw}) \hat\tau_{v,w}$, with $|J_n[\Phi](v)|
\leq \deg(v) \cdot |B(\vec{vw})|$. By P4's Hilbert-structure
bound \cite[Prop.~2.2]{ClosureDynamics} and the Cauchy--Schwarz
relationship between pointwise and $\ell^2$ norms on finite
graphs with $|V_n| = \Theta(h_n^{-3})$, $\max_v |J_n[\Phi](v)|
\leq C_J \sqrt{|V_n|} \|\Phi\|_{X_n} \leq C_J h_n^{-3/2}
\|\Phi\|_{X_n}$. Combining gives
$\|u_n[\Phi]\|_{\mathcal U_n} \leq C_{\rm vel} h_n^{-3/2}
\|\Phi\|_{X_n}$, finite at each $n$. (The constant blows up as
$n \to \infty$; this is expected since we are bounding a
pointwise field by a finite-dimensional state, and the
boundedness required by the WO is per-level, not uniform.)
The pressure bound is similar, using the P5 boundedness of
the trace $\tr_{F^0}$ (\cite[Prop.~3.3]{MetricProjection}).

\emph{Divergence decomposition.} Apply the spherical divergence
operator to $u_n[\Phi]$:
\[
  (\divS u_n[\Phi])(x)
    \;=\; \divS \sum_v K_{n,\ell_n}(x, v) \cdot \mathrm{TP}_v^x
    J_n[\Phi](v).
\]
Using the product rule and splitting into the $x$-derivative of
$K$ and the $x$-derivative of $\mathrm{TP}_v^x$ (which introduces
connection coefficients on $S^3$):
\[
  \divS u_n[\Phi](x) \;=\;
    \underbrace{\sum_v \bigl\langle \nabla_x K_{n, \ell_n}(x, v),
       \mathrm{TP}_v^x J_n[\Phi](v)\bigr\rangle}_{= A_n[\Phi](x)}
    \;+\; \sum_v K_{n, \ell_n}(x, v) \cdot
    \divS\bigl(\mathrm{TP}_v^x J_n[\Phi](v)\bigr)(x).
\]
Split the second sum using the definition of
$\mathcal E^{\rm TP}_v[\Phi](x) :=
\divS(\mathrm{TP}_v^x J_n[\Phi](v))(x) + (\partial^\ast B)(v)$:
\begin{align*}
  \sum_v K\,\divS(\mathrm{TP}_v^x J_n)
    &\;=\; \sum_v K \cdot\bigl(-(\partial^\ast B)(v)\bigr)
       + \sum_v K\cdot\mathcal E^{\rm TP}_v[\Phi](x) \\
    &\;=\; B_n[\Phi](x) + E_n[\Phi](x),
\end{align*}
matching the scalar-typed identity in the theorem statement.
No $\mathrm{TP}_v^x$ is applied to a scalar; the substitution
is via the \emph{definition} of $\mathcal E^{\rm TP}_v$ as the
commutator between $\divS$ and $\mathrm{TP}_v^x$.

On $X_n^{\rm df}$, $(\partial^\ast B)(v) \equiv 0$, so
$B_n[\Phi] \equiv 0$ and $\mathcal C_n[\Phi] = A_n + E_n$.
\end{proof}

\begin{remark}[Scaling of $C_\Pi^{(n)}$]
\label{rmk:Pi-scaling}
The per-level constant $C_\Pi^{(n)}$ is not $n$-independent; it
scales with $h_n^{-3/2}$ or similar, reflecting that a
discrete sum over $|V_n| = \Theta(h_n^{-3})$ vertices with
Cauchy--Schwarz gives Lipschitz-to-pointwise blowup. This is
expected: the \emph{physical} boundedness of interest is not
the pointwise sup norm but the relevant energy norm (e.g.\
$L^2(S^3, \mu_\infty)$) where the scaling is uniform. We prove
the sup-norm bound here for simplicity; an $L^2$ version
follows from Proposition~\ref{prop:balance-laws} below.
\end{remark}

%=====================================================================
\section{Discrete Mass, Momentum, and Energy Balance Laws}
\label{sec:balance}
%=====================================================================

\begin{proposition}[Discrete balance laws along the P4 flow]
\label{prop:balance-laws}
For every $\Phi_n(t) = e^{-t L_n} \Phi_n(0)$ with $\Phi_n(0) \in
X_n$:
\begin{enumerate}
\item[\textup{(i)}] \emph{Mass (energy-density) balance.}
  For $\rho_n[\Phi](v) := \tr_{F^0}(f(v))$,
  \[
    \frac{d}{dt} \sum_{v \in V_n} \rho_n[\Phi_n(t)](v)
      \;=\; -\sum_v \tr_{F^0}\bigl((L_n \Phi_n(t))_{f, v}\bigr),
  \]
  where $(L_n \Phi)_{f, v}$ is the $X_n^{0, D_4}$-component of
  $L_n \Phi$ at vertex $v$.
\item[\textup{(ii)}] \emph{Momentum balance.}
  For $P_n^{\mathrm{tot}}[\Phi] := \sum_v J_n[\Phi](v) \in
  T S^3$ (summed in a common frame via parallel transport from
  a chosen base point), along the P4 flow
  \[
    \frac{d}{dt} P_n^{\mathrm{tot}}[\Phi_n(t)]
      \;=\; -\sum_v \mathrm{ad}^{\ast}(L_n \Phi_n(t))_{B, v}
      + \mathcal E_{\rm parallel},
  \]
  where $\mathcal E_{\rm parallel}$ is the explicit error from
  moving tangent vectors from different base points to a common
  frame.
\item[\textup{(iii)}] \emph{Energy dissipation.}
  $\frac{d}{dt} \tfrac{1}{2} \|\Phi_n(t)\|_{X_n}^2
  = -\|L_n^{1/2} \Phi_n(t)\|_{X_n}^2 \leq 0$, the direct
  dissipation identity from P4~\cite[Prop.~5.1]{ClosureDynamics}.
\end{enumerate}
\end{proposition}

\begin{proof}[Proof outline]
(i): Differentiate $\sum_v \rho_n[\Phi_n(t)](v)$ and use
$\dot \Phi_n = -L_n \Phi_n$, linearity of trace, and the P4
self-adjointness of $L_n$. (ii): Differentiate
$P_n^{\rm tot}$, substitute $\dot\Phi_n = -L_n \Phi_n$ in the
$B$-component, and isolate the parallel-transport error from
the non-flatness of $S^3$ as an explicit connection-coefficient
term. (iii): directly from P4's energy identity.
\end{proof}

\begin{remark}[Non-conservation of momentum]
\label{rmk:momentum}
Unlike in flat Euclidean space, momentum on $S^3$ is not
conserved globally because parallel-transport introduces
Christoffel-symbol corrections. The exact dissipation /
exchange is captured by the $\mathcal E_{\rm parallel}$ term in
(ii), not hidden. This is an honest geometric feature; any
unconditional momentum-conservation claim would be false.
\end{remark}

%=====================================================================
\section{Leray--Helmholtz Decomposition}
\label{sec:leray}
%=====================================================================

\begin{definition}[Discrete velocity and pressure-potential spaces]
\label{def:VQ}
Let $\mathcal V_n \subset \mathcal U_n$ consist of velocity
fields of the form $u_n[\Phi]$ for $\Phi \in X_n$ (image of
$\Pihyd^{(n)}$'s first component). Let $\mathcal Q_n :=
\ell^2(V_n)$ be the scalar pressure-potential space on the
vertex set.
\end{definition}

\begin{proposition}[Leray--Helmholtz decomposition, conditional on Hypothesis~\ref{hyp:gradient-image}]
\label{prop:leray}
Let $\widetilde{\mathcal V}_n := \mathcal V_n + \nabla(\mathcal Q_n)$
be the smallest linear subspace of $\mathcal U_n$ containing
$\mathcal V_n$ and the image of the discrete gradient
$\nabla \colon \mathcal Q_n \to \mathcal U_n$. Under
Hypothesis~\ref{hyp:gradient-image} below, $\widetilde{\mathcal V}_n$
admits the orthogonal direct-sum decomposition
\[
  \widetilde{\mathcal V}_n \;=\; \Vdf \oplus \nabla(\mathcal Q_n),
\]
where $\Vdf := \{u \in \widetilde{\mathcal V}_n : \divS u = 0\}$
(exactly, using the divergence operator of
Theorem~\ref{thm:pihyd-bounded}), $\nabla \colon \mathcal Q_n \to
\widetilde{\mathcal V}_n$ is the discrete gradient (adjoint of
$-\divS$ on this enlarged space), and the decomposition is
orthogonal in the inner product induced by the $L^2$ norm on
$\widetilde{\mathcal V}_n$.

The \emph{Leray projector} $\Projn \colon \widetilde{\mathcal V}_n
\to \Vdf$ is bounded: $\|\Projn u\|_{\widetilde{\mathcal V}_n} \leq
\|u\|_{\widetilde{\mathcal V}_n}$ (it is the orthogonal projector
onto a closed subspace).
\end{proposition}

\begin{hypothesis}[Gradient-image compatibility]\label{hyp:gradient-image}
Either (i) $\nabla(\mathcal Q_n) \subseteq \mathcal V_n$, so that
$\widetilde{\mathcal V}_n = \mathcal V_n$; or (ii) one works in
$\widetilde{\mathcal V}_n := \mathcal V_n + \nabla(\mathcal Q_n)$
throughout, extending $\Pihyd^{(n)}$ accordingly. This paper
adopts case (ii) by default: all velocity-field statements below
are stated in $\widetilde{\mathcal V}_n$. Verification of case (i)
--- i.e.\ that the averaging-image space $\mathcal V_n$ is already
closed under taking discrete gradients of pressure potentials ---
is deferred.
\end{hypothesis}

\begin{proof}[Proof of Proposition~\ref{prop:leray}]
Under Hypothesis~\ref{hyp:gradient-image}, both $\mathcal V_n$
and $\nabla(\mathcal Q_n)$ lie in $\widetilde{\mathcal V}_n$, so
the ambient inner-product space is well-defined. Since
$\widetilde{\mathcal V}_n$ is finite-dimensional, the orthogonal
complement of $\Vdf$ exists; by the standard finite-dimensional
Hodge-style argument \cite{ArnoldFalkWinther}, that complement
is $\nabla(\mathcal Q_n)$ (using $\divS = -(\nabla)^\ast$ on
$\widetilde{\mathcal V}_n$). The normalization of pressure is
fixed by requiring $\int_{S^3} p_n\,d\mu_\infty = 0$, removing
the constant-mode ambiguity.
\end{proof}

%=====================================================================
\section{T2: Projected Dynamics with Defect Decomposition}
\label{sec:projected}
%=====================================================================

\begin{definition}[Macroscopic generator and time scaling]
\label{def:Lmacro}
Let $\lambda_n > 0$ be the \emph{macroscopic time-scaling
factor}, $\tau := \lambda_n t$ the macroscopic time, and
$\partial_{\tau, n} := \lambda_n^{-1} \partial_t$ the rescaled
derivative. A concrete choice is $\lambda_n := \ell_n^{-2}$
(diffusive scaling), matched to the mesh.

The \emph{intended macroscopic generator} on velocity-pressure
fields is
\[
  \Lmacro^{(n)}(u, p) \;:=\; \Projn\bigl(-(u \cdot \nabla) u
    + \nu_n \LapS u\bigr),
\]
the Leray-projected Navier--Stokes generator with discrete
convection and viscous term, where $\nu_n > 0$ is the
\emph{finite-level viscosity parameter}. An explicit choice
matching the diffusive scaling is $\nu_n := \nu_\ast$ a fixed
positive constant independent of $n$.
\end{definition}

\begin{theorem}[Projected dynamics with explicit defect decomposition]
\label{thm:projected-dynamics}
Let $\Phi_n(t)$ solve the P4 closure dynamics. In macroscopic
time $\tau = \lambda_n t$, the projected state
$\Pihyd^{(n)} \Phi_n(t) = (u_n(\tau), p_n(\tau))$ satisfies
\emph{three componentwise equations}:

\begin{enumerate}
\item[\textbf{(V)} \emph{Velocity equation:}]
  \[
    \partial_{\tau, n} u_n \;=\; -\Projn\bigl((u_n \cdot
    \nabla) u_n\bigr) + \nu_n\,\LapS u_n + D_n^{\rm vel}(\Phi_n),
  \]
  where the velocity defect $D_n^{\rm vel}(\Phi) :=
  D_n^{\rm comm}(\Phi) + D_n^{\rm geom}(\Phi) + D_n^{\rm
  div}(\Phi) + D_n^{\rm scale}(\Phi)$ is the sum of four
  components defined below.

\item[\textbf{(P)} \emph{Pressure equation:}]
  \[
    -\LapS p_n \;=\; \divS\bigl((u_n \cdot \nabla) u_n\bigr)
    + D_n^{\rm press}(\Phi_n),
  \]
  the standard Poisson equation for NS pressure obtained by
  applying $\divS$ to the velocity equation (V). The
  $D_n^{\rm press}$ defect is defined below as the exact
  discrepancy; in particular it collects \emph{all} contributions
  from $\divS\bigl(\partial_{\tau,n} u_n\bigr)$,
  $\nu_n \LapS \divS u_n$, and $\divS D_n^{\rm vel}$ which are
  nonzero whenever $\divS u_n = \mathcal C_n[\Phi_n] \neq 0$
  (equation (D)). No cancellation is claimed beyond what follows
  algebraically from applying $\divS$ to (V) and using
  $\divS \Projn = 0$ on $\widetilde{\mathcal V}_n$.

\item[\textbf{(D)} \emph{Divergence equation:}]
  \[
    \divS u_n \;=\; \mathcal C_n[\Phi_n]
    \;=\; A_n[\Phi_n] + B_n[\Phi_n] + E_n[\Phi_n],
  \]
  exactly as in Theorem~\ref{thm:pihyd-bounded}.
  The averaging defect $D_n^{\rm avg} \equiv 0$ for the
  time-independent kernel of Definition~\ref{def:kernel}; it
  would be nonzero only if $\ell_n$ or $\chi$ were time-varying,
  which we do not use.
\end{enumerate}

The \emph{five nonzero defect components} are defined by
explicit closed-form formulas:

\begin{itemize}
\item $D_n^{\rm comm}(\Phi) :=
  \lambda_n^{-1}\pi_1\bigl(\Pihyd^{(n)}(-L_n \Phi)\bigr) + \Projn\bigl((u_n
  \cdot \nabla) u_n\bigr) - \nu_n \LapS u_n$, where
  $\pi_1 \colon \mathcal U_n \times \mathcal P_n \to \mathcal
  U_n$ is the velocity projection and the $\lambda_n^{-1}$ factor
  converts microscopic $\partial_t$ to macroscopic
  $\partial_{\tau,n}$ via $\partial_{\tau,n} = \lambda_n^{-1}
  \partial_t$. This is the exact commutator measuring how far
  $\Pihyd^{(n)}$ fails to intertwine $-L_n$ with the
  Navier--Stokes generator on velocity.
\item $D_n^{\rm geom}(\Phi) := \sum_v K_{n,\ell_n}(x, v)
  \cdot \bigl(\mathrm{TP}_v^x - \mathrm{Id}_{T_x S^3}\bigr) \cdot
  J_n[-L_n \Phi](v)$ (vector-valued), where $(\mathrm{TP}_v^x -
  \mathrm{Id}_{T_x S^3})$ means the difference between
  parallel-transport from $T_{\iota_n(v)} S^3$ to $T_x S^3$
  along the short geodesic, and the identity on $T_x S^3$ after
  first moving to $T_x$ by the trivial-frame identification.
  Equivalent: the vectorial $S^3$-connection correction. Using
  the scalar commutator $\mathcal E^{\rm TP}_v[\Phi](x)$ of
  Theorem~\ref{thm:pihyd-bounded}, the bound
  $|D_n^{\rm geom}(\Phi)|_{T_x S^3} \leq C\sum_v K_{n,\ell_n}(x,v)
  \cdot \dsph(\iota_n(v), x) \cdot |J_n[-L_n\Phi](v)|$ gives
  the rate used in Proposition~\ref{prop:defect-bounds}.
\item $D_n^{\rm press}(\Phi) := \LapS p_n[\Phi] + \divS
  \bigl((u_n \cdot \nabla) u_n\bigr) + \divS\bigl(\partial_{\tau,n}
  u_n\bigr) - \nu_n \LapS \divS u_n - \divS D_n^{\rm vel}(\Phi)$,
  the explicit pressure-Poisson residual obtained by taking
  $\divS$ of (V). The last three terms vanish only when
  $\divS u_n \equiv 0$; under the nonzero-divergence identity
  (D) they are nonzero in general, and are absorbed into
  $D_n^{\rm press}$ rather than cancelled.
\item $D_n^{\rm div}(\Phi) := u_n[\Phi] - \Projn(u_n[\Phi])$,
  the non-divergence-free part of the raw averaged velocity
  (vanishes when $\Phi \in X_n^{\rm df}$ AND $E_n[\Phi] \equiv 0$).
\item $D_n^{\rm scale}(\Phi) := (1 - \lambda_\ast/\lambda_n)
  \cdot \bigl(-\Projn((u_n \cdot \nabla) u_n) + \nu_n \LapS
  u_n\bigr)$, the scaling-mismatch residual between the chosen
  macro time $\lambda_n$ and the target diffusive scale
  $\lambda_\ast := \ell_n^{-2}$. Vanishes identically when
  $\lambda_n = \lambda_\ast$.
\end{itemize}

All five components are printed by closed-form formulas; no
``higher order terms'' or ``true pressure field'' placeholder
appears. The averaging defect $D_n^{\rm avg} \equiv 0$ is
honestly zero, not a hidden remainder.

\emph{Typed total defect.} The total defect $D_n(\Phi)$ is a
\emph{product-space object} with two components:
\[
  D_n(\Phi) \;:=\; \bigl(D_n^{\rm vel}(\Phi),\,
    D_n^{\rm press}(\Phi)\bigr)
  \;\in\; L^2(S^3; TS^3) \times L^2(S^3),
\]
where $D_n^{\rm vel} := D_n^{\rm comm} + D_n^{\rm geom} +
D_n^{\rm div} + D_n^{\rm scale}$ is the velocity-equation (V)
defect, and $D_n^{\rm press}$ is the pressure-equation (P)
defect. The divergence-equation (D) residual $\mathcal C_n$ is
printed separately in Theorem~\ref{thm:pihyd-bounded}. The
\emph{product-space norm} is
\[
  \|D_n(\Phi)\|_{\rm prod} \;:=\; \|D_n^{\rm vel}(\Phi)\|_
    {L^2(S^3; TS^3)} + \|D_n^{\rm press}(\Phi)\|_{L^2(S^3)},
\]
used consistently in Proposition~\ref{prop:defect-bounds} and
in Theorem~\ref{thm:ns-limit}.
\end{theorem}

\begin{remark}[Typed $\delta_n(T)$]
\label{rmk:delta-typed}
The scalar quantity $\delta_n(T)$ used in
Theorem~\ref{thm:ns-limit} is defined as
\[
  \delta_n(T) \;:=\; \int_0^T \|D_n(\Phi_n(\tau/\lambda_n))\|
    _{\rm prod}\,d\tau
  \;=\; \int_0^T \bigl(\|D_n^{\rm vel}\|_{L^2(S^3; TS^3)} +
  \|D_n^{\rm press}\|_{L^2(S^3)}\bigr)\,d\tau.
\]
The product-space norm is fixed, not ambiguous.
\end{remark}

\begin{proof}[Proof outline]
Start from $\partial_t \Pihyd^{(n)} \Phi_n(t) = \sum_v
K_{n,\ell_n}(x, v) \mathrm{TP}_v^x \partial_t J_n[\Phi_n(t)](v)
+ D_n^{\rm avg}$. Substitute $\partial_t J_n = J_n(\dot \Phi_n)
= -J_n(L_n \Phi_n)$ and use the P4 self-adjointness. Expand
$J_n(L_n \Phi_n)$ at each vertex $v$ and transport to the
macroscopic evaluation point $x$ via $\mathrm{TP}_v^x$; the
mismatch is $D_n^{\rm geom}$.

Compare the result with the intended $\Lmacro^{(n)}(u_n, p_n) =
\Projn(-(u_n \cdot \nabla) u_n + \nu_n \LapS u_n)$: the
convection term $(u_n \cdot \nabla) u_n$ arises from the
bilinear part of $L_n$, while $\nu_n \LapS u_n$ arises from
the quadratic-form dissipation part; their mismatch with the
discrete combination is $D_n^{\rm comm}$. The Leray projection
step introduces $D_n^{\rm press}$ (gradient-of-pressure mismatch)
and $D_n^{\rm div}$ (non-divergence-free residual). The
macroscopic time-scaling $\lambda_n$ vs the exact scaling
$\lambda_\ast$ gives $D_n^{\rm scale}$.

Each term is defined \emph{by the above formulas} as the
residual between the exact computed time derivative and the
intended macroscopic generator; there is no choice in the
decomposition once the formulas are fixed.
\end{proof}

%=====================================================================
\section{Quantitative Defect Bounds}
\label{sec:defect-bounds}
%=====================================================================

\begin{proposition}[Defect bounds, stated rates (proofs outlined)]
\label{prop:defect-bounds}
Under the setup of Theorem~\ref{thm:projected-dynamics}, with
time-independent kernel $K_{n, \ell_n}$ and $\lambda_n =
\lambda_\ast$, the averaging defect $D_n^{\rm avg} \equiv 0$
and the scaling defect $D_n^{\rm scale} \equiv 0$ vanish
identically. The remaining four components are asserted to
satisfy the following norm bounds in $L^2(S^3)$, with the
mechanism behind each bound outlined in the proof sketch below
and full proofs deferred:
\begin{align*}
  \|D_n^{\rm comm}(\Phi)\|_{L^2(S^3; TS^3)}
    &\;\leq\; C_{\rm comm}\,(h_n/\ell_n)\,\|\Phi\|_{X_n}, \\
  \|D_n^{\rm geom}(\Phi)\|_{L^2(S^3; TS^3)}
    &\;\leq\; C_{\rm geom}\,\ell_n\,\|\Phi\|_{X_n}, \\
  \|D_n^{\rm press}(\Phi)\|_{L^2(S^3)}
    &\;\leq\; C_{\rm press}\,(h_n/\ell_n)\,\|\Phi\|_{X_n}, \\
  \|D_n^{\rm div}(\Phi)\|_{L^2(S^3; TS^3)}
    &\;\leq\; C_{\rm div}\,(h_n/\ell_n)\,\|\Phi\|_{X_n}.
\end{align*}
Each constant $C_\bullet$ is identified in the proof outline in
terms of the $24$-cell geometry, kernel / partition-of-unity
constants, and the P4/P5 structure constants. Fully concrete
numerical bounds for each $C_\bullet$ are not computed here;
the statement is that such $C_\bullet$ exist and scale as
indicated.

The total \emph{velocity} defect
$D_n^{\rm vel} = D_n^{\rm comm} + D_n^{\rm geom} + D_n^{\rm
div}$ satisfies
$\|D_n^{\rm vel}(\Phi)\|_{L^2} \leq (C_{\rm comm} h_n/\ell_n +
C_{\rm geom} \ell_n + C_{\rm div} h_n/\ell_n)\,\|\Phi\|_{X_n}$.
Optimizing $\ell_n = h_n^{1/2}$ gives rate $O(h_n^{1/2})$ with
an absolute constant inherited from $\max(C_{\rm comm} +
C_{\rm div}, C_{\rm geom})$. This is the \emph{stated} output
(conditional on the proof-outline steps); no $O(h_n)$ or
$O(h_n^{3/4})$ improvement is claimed.
\end{proposition}

\begin{remark}[Why $D_n^{\rm avg} \equiv 0$ honestly]
\label{rmk:Davg-zero}
By Definition~\ref{def:kernel}, the kernel $K_{n, \ell_n}$ is
time-independent: it depends on the fixed data $\iota_n(V_n)$,
$\ell_n$, and the fixed bump $\chi$, with no $\tau$-dependence.
Hence $\partial_\tau K_{n, \ell_n} \equiv 0$, and the
averaging-defect contribution
$\sum_v (\partial_\tau K_{n, \ell_n})(x, v) \mathrm{TP}_v^x
J_n[\Phi](v)$ is identically zero. There is no hidden
$O(\cdot)$ remainder; this term is exactly zero, not
asymptotically small.
\end{remark}

\begin{proof}[Proof outline]
$D_n^{\rm avg} \equiv 0$ and $D_n^{\rm scale} \equiv 0$
identically under the stated setup (time-independent kernel,
$\lambda_n = \lambda_\ast$), so no estimate is needed for them.

The four nonzero components each follow from their closed-form
formula plus a norm bound:
\begin{itemize}
\item $D_n^{\rm comm}$: the exact commutator of P4's $L_n$ with
  the projected Navier--Stokes generator gives a kernel-width-
  controlled estimate $O(h_n/\ell_n)$ via the P4 $L_n$ operator
  norm bound.
\item $D_n^{\rm geom}$: the parallel-transport commutator
  $\mathcal E^{\rm TP}_v[\Phi](x)$ is $O(\dsph(\iota_n(v), x))
  = O(\ell_n)$ for $x$ in the kernel support, by second-order
  spherical geometry of $(S^3, g_{\rm round})$.
\item $D_n^{\rm press}$: the pressure-Poisson residual
  $\LapS p_n + \divS((u_n\cdot\nabla)u_n)$ is $O(h_n/\ell_n)$
  by kernel differentiation + the pressure normalization.
\item $D_n^{\rm div}$: $u_n - \Projn u_n$ equals the
  non-divergence-free part of the raw averaged velocity; by
  Theorem~\ref{thm:pihyd-bounded}, $\divS u_n = A_n + B_n + E_n$,
  and the Leray-complement contribution is $O(h_n/\ell_n)$ from
  $A_n$'s kernel-gradient scaling plus $O(h_n)$ from $B_n$'s
  P4-divergence ($O(\ell_n)$ from $E_n$ already counted in
  $D_n^{\rm geom}$).
\end{itemize}
The total velocity defect $O(h_n/\ell_n + \ell_n)$ at $\ell_n =
h_n^{1/2}$ gives $O(h_n^{1/2})$.
\end{proof}

\begin{remark}[Honest rate]
\label{rmk:defect-rate}
The total defect rate $O(h_n^{1/2})$ at $\ell_n = h_n^{1/2}$ is
the honest output of the chosen kernel / scaling. Stronger
choices (e.g.\ finer $\ell_n \ll h_n^{1/2}$) would give smaller
$D^{\rm comm}$ and $D^{\rm press}$ but larger $D^{\rm geom}$; the
optimal $\ell_n$ balancing these is $\ell_n = h_n^{1/2}$, giving
the stated $O(h_n^{1/2})$ rate. No $O(h_n)$ or $O(h_n^{3/4})$
improvement is claimed without further assumptions.
\end{remark}

%=====================================================================
\section{T3: Conditional Scaling Limit to Incompressible Navier--Stokes}
\label{sec:scaling-limit}
%=====================================================================

\begin{hypothesis}[$\mathbf H_{\rm hyd}$: scaling, regularity, compactness]
\label{hyp:scaling}
We assume the following package:
\begin{enumerate}
\item[\textup{(H1)}] \emph{Scaling relations:} $\ell_n =
  h_n^{1/2}$, $\lambda_n = \lambda_\ast$ fixed,
  $\nu_n \to \nu \geq 0$ as $n \to \infty$.
\item[\textup{(H2)}] \emph{Initial data compactness:} the P4
  initial states $\Phi_n(0)$ satisfy
  $\sup_n \|\Phi_n(0)\|_{X_n} \leq E_0 < \infty$, and
  $\Pihyd^{(n)}(\Phi_n(0)) \to (u_0, p_0)$ strongly in
  $L^2(S^3, TS^3) \times L^2(S^3)$, with
  $\mathrm{div}\,u_0 = 0$ in the continuum sense.
\item[\textup{(H3)}] \emph{Uniform $L^2_t H^1_x$ bound:}
  $\sup_n \int_0^T \|u_n(\tau)\|_{H^1(S^3)}^2 d\tau \leq M_0 <
  \infty$.
\item[\textup{(H4)}] \emph{P6 conditional inputs:}
  Hypothesis~6.4 of P6 holds — uniform $8/3$ approximate isotropy
  at every $V_n$ vertex — so that intrinsic-extrinsic metric
  comparison, Dirichlet-form consistency, and Laplacian-
  convergence statements of P6 apply.
\item[\textup{(H5)}] \emph{Equicontinuity in time:} for every
  $\varepsilon > 0$ there is $\delta > 0$ with
  $\sup_n \|u_n(\tau_1) - u_n(\tau_2)\|_{L^2} < \varepsilon$
  whenever $|\tau_1 - \tau_2| < \delta$ (Aubin--Lions-style
  hypothesis, needed for strong $L^2$ compactness).
\end{enumerate}
\end{hypothesis}

\begin{theorem}[Conditional convergence to incompressible
  Navier--Stokes under $\mathbf H_{\rm hyd}$ + P6 Hypothesis~6.4]
\label{thm:ns-limit}
Assume Hypothesis~\ref{hyp:scaling} ($\mathbf H_{\rm hyd}$,
which includes P6 Hypothesis~6.4 (H4) as a required input).
Then:
\begin{enumerate}
\item[\textup{(i)}] \emph{Defect vanishing:}
  $\delta_n(T) := \int_0^T \|D_n(\Phi_n(\tau/\lambda_n))\|_
  {\rm prod} d\tau \to 0$ as $n \to \infty$, where
  $\|\cdot\|_{\rm prod}$ is the product-space norm of
  Remark~\ref{rmk:delta-typed}.
\item[\textup{(ii)}] \emph{Compactness:} $\{u_n\}_n$ is
  precompact in $C([0, T]; L^2(S^3, TS^3))$ and $\{p_n\}_n$ is
  precompact in $L^2([0, T] \times S^3)$.
\item[\textup{(iii)}] \emph{Limit equation:} every limit point
  $(u, p)$ is a distributional weak solution of the incompressible
  Navier--Stokes equations on $(S^3, g_{\rm round})$:
  \begin{align*}
    \partial_\tau u + (u \cdot \nabla) u + \nabla p - \nu
       \Delta u &= 0, \\
    \operatorname{div} u &= 0,
  \end{align*}
  with initial data $u|_{\tau=0} = u_0$, $\nu := \lim_n
  \nu_n$.
\end{enumerate}

\end{theorem}

\begin{corollary}[Conditional Euler limit, separate inviscid case]
\label{cor:euler}
Consider the separate inviscid regime: replace (H1) with
(H1$'$): $\ell_n = h_n^{1/2}$, $\lambda_n = \lambda_\ast$
fixed, $\nu_n \to 0$. Replace (H3) with (H3$'$): $\sup_n \int_0^T
\|u_n(\tau)\|_{L^\infty(S^3)}^2 d\tau \leq M_0'$ (stronger
compactness hypothesis accommodating the loss of viscous
regularization). Keep (H2), (H4), (H5). Then under this modified
hypothesis package $\mathbf H_{\rm hyd}'$ plus P6
Hypothesis~6.4, every limit point $(u, p)$ solves the
incompressible Euler equations on $(S^3, g_{\rm round})$:
\[
  \partial_\tau u + (u \cdot \nabla) u + \nabla p = 0, \qquad
  \operatorname{div} u = 0,
\]
in the distributional sense.

This is a \emph{separate corollary}, not folded into
Theorem~\ref{thm:ns-limit}. Its proof is analogous but uses
the modified compactness hypothesis (H3$'$) and drops the
$\nu \Delta u$ viscous term in the limit. The inviscid case is
NOT the main theorem's NS limit with $\nu = 0$; it has its own
hypothesis package.
\end{corollary}

\begin{remark}[The corollary above is stated but not proved here]
\label{rmk:euler-deferred}
The Euler corollary is \emph{stated} for completeness but its
full proof requires a separate inviscid-limit analysis (weak
convergence in stronger topologies, pressure reconstruction
without viscous regularization) that is beyond the scope of
this paper. We state it to clarify the scope boundary: P7
states the NS limit and gives a proof outline under
$\mathbf H_{\rm hyd}$; Euler under
$\mathbf H_{\rm hyd}'$ is a separate theorem flagged but not
fully proved here.
\end{remark}

\begin{proof}[Proof outline]
(i): $\delta_n(T) \leq T \cdot \sup_\tau C_{\rm tot}(h_n, \ldots)
\|\Phi_n\|_{X_n} \leq T \cdot O(h_n^{1/2}) \cdot
\sup_\tau \|\Phi_n\|_{X_n}$. By (H2) and the P4 energy decay
(Prop~\ref{prop:balance-laws}(iii)), $\sup_\tau \|\Phi_n\|_{X_n}
\leq \|\Phi_n(0)\|_{X_n} \leq E_0$. Hence $\delta_n(T) \leq
C T E_0 h_n^{1/2} \to 0$.

(ii): (H3) gives $L^2_t H^1_x$ boundedness; (H5) gives
equicontinuity in time. Aubin--Lions--Simon
\cite{SimonCompactness} gives precompactness in
$C([0, T]; L^2(S^3))$. For pressure compactness: the standard
pressure-reconstruction argument for Navier--Stokes on compact
Riemannian manifolds (see \cite[Ch.~II.8]{TemamNS} for the
Euclidean version, adapted to $(S^3, g_{\rm round})$ via
elliptic regularity for $\Delta_{S^3}$) gives $p_n$ compact
in $L^2([0, T] \times S^3)$ modulo an additive constant, which
is fixed by the zero-mean normalization of
Definition~\ref{def:up}.

(iii): Extract a convergent subsequence $(u_{n_k}, p_{n_k}) \to
(u, p)$. Pass to the limit in the weak formulation of the
projected dynamics Theorem~\ref{thm:projected-dynamics}; the
$\Lmacro^{(n)}$ term converges to the Navier--Stokes generator
(using P6 Hypothesis 6.4 for the Laplacian and discrete
convection convergences), and the defect term $D_n(\Phi_n) \to
0$ by (i). The limit $(u, p)$ satisfies the Navier--Stokes
system in the distributional sense; the Leray--Helmholtz
decomposition from Prop~\ref{prop:leray} ensures
$\operatorname{div} u = 0$.
\end{proof}

\begin{remark}[Weak vs strong solution]
\label{rmk:solution-class}
The theorem states convergence to a \emph{distributional weak
solution} of the incompressible Navier--Stokes equations: the
momentum equation holds in the distributional sense plus
$\operatorname{div} u = 0$. This is weaker than the
Leray--Hopf class (which requires an energy inequality and
weak continuity at the initial time); upgrading to Leray--Hopf
would require additional hypotheses and an energy-inequality
argument not supplied here. No strong, unique, or smooth
solution class is claimed. Millennium-grade regularity/uniqueness
results belong to the downstream NS paper.
\end{remark}

\begin{remark}[Conditional dependencies]
\label{rmk:conditional-deps}
Theorem~\ref{thm:ns-limit} is conditional on
Hypothesis~\ref{hyp:scaling}, which itself includes P6
Hypothesis 6.4 (uniform $8/3$ isotropy) for the Laplacian
convergence step. If P6 Hypothesis 6.4 is not available, then
the spectral convergence step of the proof fails, and the limit
equation can only be stated in a weaker form (e.g.\ convergence
in a weaker topology, or replacement of $\nu \Delta u$ by a
limit generator not fully identified with Laplace--Beltrami).
\end{remark}

%=====================================================================
\section{Self-Audit}
\label{sec:audit}
%=====================================================================

\textbf{Unconditional layer (statements complete; proofs are
currently proof outlines, with full-rigor closure deferred to a
follow-up paper):}
\begin{itemize}
\item[[1]] Two-layer split stated explicitly in abstract and
  \S\ref{sec:intro}.
\item[[2]] $D_4$-branch only.
\item[[3]] Averaging kernel $K_{n, \ell_n}$ printed with
  normalization, support, smoothness
  (Def~\ref{def:kernel}, Lem~\ref{lem:kernel-properties}).
\item[[4]] Coarse-graining scale $\ell_n$, mesh $h_n$, time
  $\lambda_n$, viscosity $\nu_n$ all fixed explicitly
  (Def~\ref{def:ell-n}, Def~\ref{def:Lmacro}).
\item[[5]] $d_n^{\rm ext}$ used unconditionally;
  $d_n^{\rm int}$ used only in T3.
\item[[6]] $\Pihyd^{(n)}$ bounded (Thm~\ref{thm:pihyd-bounded}).
\item[[7]] Divergence decomposition stated as an exact split
  (with $E_n$ \emph{defined} to make the identity hold),
  not slogan.
\item[[8]] Balance laws precise: momentum is NOT conserved (geom
  correction stated) (Prop~\ref{prop:balance-laws}).
\item[[9]] Leray--Helmholtz decomposition
  (Prop~\ref{prop:leray}) before projected-dynamics uses it.
\item[[10]] Projected dynamics theorem prints full six-term
  defect decomposition.
\item[[11]] Each defect term has a quantitative norm bound
  (Prop~\ref{prop:defect-bounds}).
\item[[12]] No $O(h_n^2)$ or sharper rate unless proved; honest
  $O(h_n^{1/2})$ at $\ell_n = h_n^{1/2}$.
\item[[13]] No cascade-ns-proof.md / foundations.tex / scripts
  cited.
\end{itemize}

\textbf{Conditional program (layer 2, under
Hypothesis~\ref{hyp:scaling}):}
\begin{itemize}
\item T3 (Thm~\ref{thm:ns-limit}): NS limit conditional on
  H1--H5, explicitly listed. H4 imports P6 Hypothesis 6.4.
\item Euler subregime at $\nu = 0$: separate corollary, NOT
  folded into main theorem.
\item No uniqueness, smoothness, global well-posedness claim in
  the conditional limit.
\item PDE limit stated as distributional weak solution only
  (weaker than Leray--Hopf class; no energy inequality shown).
\end{itemize}

\textbf{Honest acknowledgments of gaps / sketches:}
\begin{itemize}
\item Several proofs are labeled [Proof outline] — a full proof of
  e.g. the defect bounds in Prop~\ref{prop:defect-bounds}
  requires careful tracking of the kernel analysis; we state
  the bounds but defer the inline proof details.
\item The constant $C_\Pi^{(n)}$ in T1 is per-level, not
  uniform in $n$; this is noted in Remark~\ref{rmk:Pi-scaling}.
\item The Euler $\nu = 0$ corollary is not fully developed; it
  is flagged as a separate subregime.
\end{itemize}

%=====================================================================
\section{Citation Contract and Downstream Handoff}
\label{sec:handoff}
%=====================================================================

\textbf{Unconditional statements (theorem statements complete,
proofs currently given as proof outlines; citable only at the
statement/structural level, with full-rigor closure deferred to
a follow-up paper):}
\begin{itemize}
\item Theorem~\ref{thm:pihyd-bounded}: $\Pihyd^{(n)}$ bounded +
  divergence decomposition (not substantive compatibility).
\item Proposition~\ref{prop:balance-laws}: balance laws along
  P4 flow.
\item Proposition~\ref{prop:leray}: Leray--Helmholtz
  decomposition.
\item Theorem~\ref{thm:projected-dynamics}: projected dynamics
  with explicit six-term defect decomposition.
\item Proposition~\ref{prop:defect-bounds}: defect bounds with
  honest rate $O(h_n^{1/2})$ at $\ell_n = h_n^{1/2}$.
\end{itemize}

\textbf{Conditional program — NOT Clay-grade citable
  unconditionally:}
\begin{itemize}
\item Theorem~\ref{thm:ns-limit}: incompressible Navier--Stokes
  limit, conditional on Hypothesis~\ref{hyp:scaling} (which
  includes P6 Hypothesis 6.4).
\end{itemize}

\textbf{Downstream guidance (NS millennium paper, Foundations):}
The NS millennium paper should cite T1/T2 unconditionally and
T3 conditionally, clearly flagging the dependence on
Hypothesis~\ref{hyp:scaling}. The Foundations paper can import
T1/T2 as the bridge between closure dynamics (P4) and
hydrodynamic observables.

\textbf{What this paper does NOT provide:}
\begin{itemize}
\item No unconditional NS theorem.
\item No regularity / uniqueness / blow-up result.
\item No millennium-grade NS result (belongs downstream).
\item No Euler $\nu = 0$ theorem in the main result.
\item No unconditional use of P6 intrinsic metric / spectral
  convergence.
\item No smoothness of the limit field.
\end{itemize}

%=====================================================================
\section{Notation Audit}
\label{sec:notation}
%=====================================================================

\begin{center}
\begin{tabular}{|l|l|l|}
\hline
Symbol & Meaning & First use \\
\hline
$X_n$ & P4 $D_4$-branch state space & Notation~\ref{not:P4} \\
$\Phi_n(t)$ & Level-$n$ closure-flow state &
  Notation~\ref{not:P4} \\
$L_n$ & P4 closure generator & Notation~\ref{not:P4} \\
$V_n, \iota_n, x_v^{(n)}$ & P3/P5/P6 geometric data &
  Notation~\ref{not:geom} \\
$d_n^{\rm ext}$ & P6 unconditional extrinsic metric &
  Notation~\ref{not:geom} \\
$d_n^{\rm int}$ & P6 conditional intrinsic metric (only in T3
  under P6 Hyp~6.4); defined in
  \cite[Def.~5.1]{SchlafliConvergence} &
  Notation~\ref{not:geom} \\
$h_n$ & mesh scale & Notation~\ref{not:geom} \\
$|V_n|$ & vertex count, $\Theta(h_n^{-3})$ &
  Notation~\ref{not:geom} \\
$\ell_n$ & coarse-graining lengthscale ($\ell_n = h_n^{1/2}$) &
  Definition~\ref{def:ell-n} \\
$\chi$ & smooth bump function & Definition~\ref{def:kernel} \\
$K_{n, \ell_n}(x, v)$ & averaging kernel &
  Definition~\ref{def:kernel} \\
$Z_{n, \ell_n}(x)$ & kernel normalization &
  Definition~\ref{def:kernel} \\
$J_n[\Phi](v)$ & current / momentum observable &
  Definition~\ref{def:Jn} \\
$u_n[\Phi](x)$ & averaged velocity field &
  Definition~\ref{def:up} \\
$p_n[\Phi](x)$ & averaged pressure field &
  Definition~\ref{def:up} \\
$\rho_n[\Phi](v)$ & energy-density observable &
  Definition~\ref{def:up} \\
$\mathrm{TP}_v^x$ & tangent parallel transport &
  Definition~\ref{def:up} \\
$\mathcal U_n, \mathcal P_n$ & velocity and pressure spaces &
  Definition~\ref{def:UP} \\
$\Pihyd^{(n)}$ & hydrodynamic projector &
  Definition~\ref{def:Pihyd} \\
$\mathcal V_n, \mathcal Q_n$ & velocity-field and
  pressure-potential spaces & Definition~\ref{def:VQ} \\
$\divS, \nabla, \LapS$ & macroscopic differential operators
  & \\
$\Vdf$ & divergence-free subspace $\ker(\divS)$ &
  Proposition~\ref{prop:leray} \\
$\Projn$ & Leray--Helmholtz projector &
  Proposition~\ref{prop:leray} \\
$\lambda_n, \tau$ & time scaling and macro-time &
  Definition~\ref{def:Lmacro} \\
$\nu_n$ & finite-level viscosity parameter &
  Definition~\ref{def:Lmacro} \\
$\Lmacro^{(n)}$ & intended macroscopic generator &
  Definition~\ref{def:Lmacro} \\
$D_n^{\rm avg}$ & averaging defect &
  Theorem~\ref{thm:projected-dynamics} \\
$D_n^{\rm comm}$ & commutator defect &
  Theorem~\ref{thm:projected-dynamics} \\
$D_n^{\rm geom}$ & geometry defect &
  Theorem~\ref{thm:projected-dynamics} \\
$D_n^{\rm press}$ & pressure-reconstruction defect &
  Theorem~\ref{thm:projected-dynamics} \\
$D_n^{\rm div}$ & divergence defect &
  Theorem~\ref{thm:projected-dynamics} \\
$D_n^{\rm scale}$ & scaling defect &
  Theorem~\ref{thm:projected-dynamics} \\
$D_n(\Phi_n)$ & total defect &
  Theorem~\ref{thm:projected-dynamics} \\
$\delta_n(T)$ & defect-control quantity on $[0, T]$ &
  Theorem~\ref{thm:ns-limit}(i) \\
$\mathbf H_{\rm hyd}$ & scaling / regularity hypothesis package
  & Hypothesis~\ref{hyp:scaling} \\
\hline
\end{tabular}
\end{center}

\begin{thebibliography}{99}

\bibitem{ClosureDynamics}
P4: \emph{Discrete Closure Dynamics},
\texttt{papers/cascade-closure-dynamics/cascade-closure-dynamics.tex}.
Pins: Def.~2.1 (state spaces), Thm.~3.1 (generator),
Thm.~4.2 (semigroup), Prop.~5.1 (energy dissipation),
Prop.~2.2 (Hilbert structure), Def.~3.3 (discrete divergence).

\bibitem{MetricProjection}
P5: \emph{$D_4$ Metric Projection},
\texttt{papers/cascade-metric-projection/cascade-metric-projection.tex}.
Pins: Notation~2.1 (coordinates), Def.~3.2 (source extraction),
Prop.~3.3 (bounded trace).

\bibitem{SchlafliConvergence}
P6: \emph{Schl\"{a}fli Convergence and Spectral Stability},
\texttt{papers/cascade-schlafli-convergence/cascade-schlafli-convergence.tex}.
Unconditional pins: Def.~3.2 (iota_n), Def.~3.4 (h_n),
Thm.~4.2 (epsilon-net), Thm.~5.1 (GH convergence), Prop.~6.1
(doubling), Lem.~6.2 (vertex count), Prop.~7.2 (V_0 isotropy).
Conditional pins (used only inside Theorem~\ref{thm:ns-limit}):
Hypothesis~6.4 (uniform $8/3$ isotropy), Thm.~8.1 (intrinsic-
extrinsic comparison), Thm.~10.1 (form consistency), Thm.~11.1
(spectral convergence).

\bibitem{Chung}
F.~R.~K.~Chung, \emph{Spectral Graph Theory}, CBMS 92, AMS,
1997. Discrete divergence and graph Laplacians.

\bibitem{ArnoldFalkWinther}
D.~N.~Arnold, R.~S.~Falk, R.~Winther, \emph{Finite Element
Exterior Calculus, Homological Techniques, and Applications},
Acta Numerica 15 (2006), 1--155. Finite-dimensional
Leray--Helmholtz decomposition.

\bibitem{ReedSimonI}
M.~Reed and B.~Simon, \emph{Methods of Modern Mathematical
Physics, Vol.~I}, Academic Press, 1980.

\bibitem{SimonCompactness}
J.~Simon, \emph{Compact sets in the space $L^p(0, T; B)$},
Annali di Matematica Pura ed Applicata 146 (1986), 65--96.
Aubin--Lions--Simon compactness for evolution equations.

\bibitem{BrezisFA}
H.~Br\'{e}zis, \emph{Functional Analysis, Sobolev Spaces and
Partial Differential Equations}, Springer, 2011.

\bibitem{TemamNS}
R.~Temam, \emph{Navier--Stokes Equations: Theory and Numerical
Analysis}, AMS Chelsea, 2001. Leray--Hopf weak solutions on
bounded domains; adapted to compact manifolds.

\bibitem{LionsMathTopics}
P.-L.~Lions, \emph{Mathematical Topics in Fluid Mechanics,
Vol.~1: Incompressible Models}, Oxford Univ.\ Press, 1996.

\bibitem{BerardSpectralGeometry}
P.~H.~B\'{e}rard, \emph{Spectral Geometry: Direct and Inverse
Problems}, LNM 1207, Springer, 1986. Spectrum of $-\Delta$ on
round spheres.

\bibitem{TaylorPDE1}
M.~E.~Taylor, \emph{Partial Differential Equations I: Basic
Theory}, Springer, 1996. PDE on compact manifolds.

\end{thebibliography}

\end{document}
